\section{Additional Results During Decoding}
\label{sec:decoding-appendix}
\raggedbottom

We extend the direct-channel and side-channel experiments to generated tokens during decoding, using Qwen2.5-0.5B and Llama-3.2-1B-Instruct.
The classification label is the generated token consumed by the measured transformer execution, which predicts the next token.
We use the Intel Core i7-10700 and AMD Ryzen 7 8700G platforms from Section~\ref{sec:architecture-results}.
All computations in this appendix run on the CPU, with INT8 matrix multiplication and FP16 outlier gathering and correction.

We use the 99.9\% profiling baseline described in Section~\ref{sec:background:quantization}.
For each projection input, we compute one threshold from the absolute FP16 activation values in the first 16 training-prompt prefills of the unquantized model.
These thresholds are fixed before quantized inference and shared across both CPUs.
At runtime, a column is selected when its maximum absolute value across token rows reaches the threshold.
We report this specified percentile setting separately from the top-k and robust-$z$ results in Section~\ref{sec:direct-results}.

\subsection{Direct-Channel Classification}
\label{sec:decoding-direct}

We use 200 encyclopedia prompts with different topics: 128 for training, 24 for validation, and 48 for testing.
Both models use greedy decoding; Qwen uses plain prompts and Llama its official chat template.
We record 32 generated tokens per response as they enter decoding, giving 4,096 training, 768 validation, and 1,536 test samples per model.
Features record selected columns in all FFN \texttt{down\_proj} layers, grouped at widths $g=1,4,16,32$.
We train a classifier with one hidden layer of 1,024 ReLU units using AdamW for 40 epochs, with learning rate $10^{-3}$ and batch size 128.
Feature selection uses training data only; validation accuracy selects the saved classifier before testing.
Tokens absent from the training vocabulary count as errors.

\begin{table}[H]
    \centering\footnotesize\setlength{\tabcolsep}{3pt}
    \caption{Direct-channel classification of 1,536 decoding tokens per model. Group size is measured in columns. Balanced accuracy averages recall over test token types.}
    \label{tab:decoding-direct}
    \begin{tabular}{@{}lrrrr@{}}
        \hline
        Model & Group & Correct & Acc. (\%) & Balanced (\%) \\\hline
        Qwen2.5-0.5B & 1 & 1156/1536 & 75.26 & 42.67 \\
        Qwen2.5-0.5B & 4 & 1146/1536 & 74.61 & 41.45 \\
        Qwen2.5-0.5B & 16 & 1151/1536 & 74.93 & 42.21 \\
        Qwen2.5-0.5B & 32 & 1151/1536 & 74.93 & 42.41 \\
        Llama-3.2-1B & 1 & 1021/1536 & 66.47 & 36.30 \\
        Llama-3.2-1B & 4 & 1020/1536 & 66.41 & 36.35 \\
        Llama-3.2-1B & 16 & 1018/1536 & 66.28 & 35.56 \\
        Llama-3.2-1B & 32 & 1012/1536 & 65.89 & 35.15 \\
        \hline
    \end{tabular}
\end{table}

Table~\ref{tab:decoding-direct} shows that grouping preserves most of the token-classification accuracy.
At $g=32$, accuracy is 74.93\% for Qwen and 65.89\% for Llama.
Predicting the most frequent training token gives 4.62\% and 6.58\%; predicting the most frequent token at each response position gives 7.88\% and 9.18\%.
The training sets contain 1,203 and 1,421 token types, covering 80.60\% and 73.96\% of the test tokens, respectively.
Neither model recovers a complete 32-token test prefix.

\subsection{Side-Channel Fidelity}
\label{sec:decoding-fidelity}

We measure the prefill and first eight decoding steps of 16 test responses, yielding 128 decoding tokens per model and platform.
The original Flush+Reload timer and scan code surround the actual FP16 gather within the model process; gathered values feed the outlier correction.
Every measured response produces the same logits and generated tokens with and without the probes.
A separate cache-hit calibration uses 1,000 samples to select a timing threshold and 1,000 to check it on each model and platform.

Decoding provides one activation row, with 32 FP16 columns per cache line.
We test one full scan with one gather, five full scans with five gathers, and stride-5 scans with 25 gathers.
The last setting probes every fifth cache line from each of five starting positions and repeats this process five times.
Thus, it changes both the scan method and the number of repeated gathers.
All three conditions are tested on each activation; the one-gather result counts one execution within that condition.
A group is positive when at least 60\% of its repeated probes indicate a hit, with votes from at least half the observed rows.
Filtering keeps the first group in each consecutive run of positive groups.

\begin{table}[H]
    \centering\footnotesize\setlength{\tabcolsep}{3pt}
    \caption{Decoding side-channel fidelity with one observed row and $g=32$. All rates are percentages. FPR precedes consensus and filtering; precision and recall follow filtering. One or five gathers use full scans; 25 gathers use repeated stride-5 scans.}
    \label{tab:decoding-cache}
    \begin{tabular}{@{}llrrrr@{}}
        \hline
        CPU & Model & Gathers & Raw FPR & Precision & Recall \\\hline
        Intel & Qwen & 1 & 68.64 & 2.96 & 20.93 \\
        Intel & Qwen & 5 & 68.51 & 3.16 & 15.67 \\
        Intel & Qwen & 25 & 9.08 & 60.07 & 55.62 \\
        Intel & Llama & 1 & 70.15 & 4.03 & 23.46 \\
        Intel & Llama & 5 & 70.01 & 4.05 & 15.57 \\
        Intel & Llama & 25 & 10.28 & 69.54 & 66.18 \\
        AMD & Qwen & 1 & 43.32 & 7.95 & 41.54 \\
        AMD & Qwen & 5 & 32.81 & 18.18 & 71.94 \\
        AMD & Qwen & 25 & 2.51 & 98.63 & 94.00 \\
        AMD & Llama & 1 & 29.88 & 15.05 & 58.19 \\
        AMD & Llama & 5 & 37.33 & 16.95 & 66.03 \\
        AMD & Llama & 25 & 7.85 & 91.42 & 93.59 \\
        \hline
    \end{tabular}
\end{table}

Table~\ref{tab:decoding-cache} reports column-group recovery.
Table~\ref{tab:decoding-classification} applies the trained $g=32$ classifier directly to these measured groups without retraining, using the same 128 tokens for its direct and side-channel results.
With 25 gathers, side-channel accuracy is 0.78--3.12 percentage points below the direct-channel result on the same tokens.
For AMD/Qwen, the 0.78-point loss has a 95\% interval of 0--2.34 points when resampling whole responses; the accuracy loss is not always below one percentage point.
One-gather accuracy is much lower, ranging from 3.91\% to 33.59\%.
No complete eight-token test prefix is recovered from the cache traces.
Small differences in CPU arithmetic change some Llama outputs across platforms, so each platform is scored against its own generated tokens.

\begin{table}[H]
    \centering\footnotesize\setlength{\tabcolsep}{4pt}
    \caption{Token classification accuracy (\%) on the same 128 decoding tokens per model and platform. Side-channel columns use filtered groups and the stated number of gathers.}
    \label{tab:decoding-classification}
    \begin{tabular}{llrrrr}
        \hline
        & & & \multicolumn{3}{c}{Side-Channel Acc.} \\
        CPU & Model & Direct & 1 & 5 & 25 \\\hline
        Intel & Qwen & 66.41 & 3.91 & 6.25 & 63.28 \\
        Intel & Llama & 60.94 & 6.25 & 3.91 & 59.38 \\
        AMD & Qwen & 66.41 & 14.06 & 39.06 & 65.62 \\
        AMD & Llama & 63.28 & 33.59 & 36.72 & 61.72 \\
        \hline
    \end{tabular}
\end{table}

\begin{table}[H]
    \centering\footnotesize\setlength{\tabcolsep}{4pt}
    \caption{FPR (\%) with gather disabled, one observed row, and ten measurement blocks. Full scan uses one scan; stride-5 uses five repetitions. Fresh flush flushes the target before each individual probe.}
    \label{tab:decoding-no-gather}
    \begin{tabular}{lrrrr}
        \hline
        CPU & Columns & Full scan & Stride-5 & Fresh flush \\\hline
        Intel & 4,864 & 67.11 & 3.24 & 0.00 \\
        Intel & 8,192 & 67.73 & 2.99 & 0.04 \\
        AMD & 4,864 & 32.96 & 14.50 & 0.00 \\
        AMD & 8,192 & 40.94 & 10.02 & 0.00 \\
        \hline
    \end{tabular}
\end{table}

Table~\ref{tab:decoding-no-gather} shows high false-positive rates even when no gather occurs.
Flushing each target again before probing reduces the FPR to 0--0.04\%, indicating that the scan itself contributes to false positives.
With four fixed selected columns, changing the input from one to 64 rows gives filtered recall of 92.5--100\% under repeated stride-5 scans on the two CPUs.
The error therefore cannot be attributed to decoding alone.

\begin{table}[H]
    \centering\footnotesize\setlength{\tabcolsep}{3pt}
    \caption{Prefill fidelity under the same percentile thresholds, one observed row, and 25 gathers. All rates are percentages; precision and recall follow filtering. The corresponding decoding results are in Table~\ref{tab:decoding-cache}.}
    \label{tab:decoding-prefill}
    \begin{tabular}{llrrr}
        \hline
        CPU & Model & Raw FPR & Precision & Recall \\\hline
        Intel & Qwen2.5-0.5B & 94.24 & 62.93 & 1.77 \\
        Intel & Llama-3.2-1B & 97.19 & 71.35 & 0.83 \\
        AMD & Qwen2.5-0.5B & 91.65 & 86.10 & 3.14 \\
        AMD & Llama-3.2-1B & 93.51 & 92.53 & 2.16 \\
        \hline
    \end{tabular}
\end{table}

Real prefill accesses are much denser under this threshold setting, and consecutive-group filtering removes many true groups (Table~\ref{tab:decoding-prefill}).
Observing 4--32 prefill rows still gives only 0.47--2.03\% filtered recall with 25 gathers.
These results use the percentile setting specified above and do not repeat the sparse setting in Table~\ref{tab:recovery}.
Overall, decoding outlier accesses contain token information, but measured classification depends on the scan method and repeated execution.
These synchronized, in-process measurements do not demonstrate recovery from an independently running victim.
